\section{Đặt vấn đề và tính cấp thiết}

\subsection{Sự hồi sinh mạnh mẽ của xu hướng nhiếp ảnh phim}
Trong hơn hai thập kỷ qua, cuộc cách mạng kỹ thuật số đã định hình lại hoàn toàn ngành công nghiệp ghi hình thị giác. Các dòng máy ảnh số hiện đại (DSLR, Mirrorless) cùng cảm biến tích hợp trên điện thoại thông minh ngày càng sở hữu độ phân giải hàng trăm megapixel, tốc độ lấy nét theo pha trong tích tắc và khả năng chụp liên tiếp hàng chục khung hình mỗi giây. Tuy nhiên, cùng với sự bão hòa của hình ảnh kỹ thuật số mang tính "mì ăn liền", một bộ phận đông đảo giới trẻ thuộc thế hệ Gen Z, các nghệ sĩ nhiếp ảnh độc lập và giới sáng tạo tại Việt Nam cũng như trên toàn thế giới đang tìm về với \textbf{Nhiếp ảnh phim truyền thống}~\cite{edwards2023tactility}.

Sự lôi cuốn của nhiếp ảnh phim đến từ nhiều yếu tố cốt lõi:
\begin{itemize}[leftmargin=2em]
    \item \textbf{Đặc tính thẩm mỹ độc bản:} Hạt phim (film grain) được tạo nên từ tinh thể bạc halogenua (Silver Halide), mang lại dải chuyển màu mượt mà, độ sâu trường ảnh hữu cơ và tông màu đặc trưng không thể tái tạo hoàn hảo bằng các bộ lọc kỹ thuật số thông thường~\cite{edwards2023tactility, su2024deep}.
    \item \textbf{Tính chiêm nghiệm và kỷ luật tư duy thị giác:} Với mỗi cuộn phim chỉ giới hạn từ 12, 24 đến 36 kiểu ảnh (đối với phim 135) hoặc 8 đến 16 kiểu (đối với phim medium format 120), người chụp buộc phải tính toán tỉ mỉ từ bố cục, ánh sáng, góc chụp cho đến khẩu độ và tốc độ màn trập trước khi bấm máy. Sự chậm rãi này rèn luyện tư duy làm việc nghiêm túc, sâu sắc cho người nghệ sĩ.
    \item \textbf{Quy trình hóa học thủ công đầy tính ma thuật:} Cảm giác tự tay pha chế hóa chất, xoay cuộn phim trong tank tráng trong buồng tối hoàn toàn, rồi hồi hộp ngắm nhìn từng dải phim âm bản lộ diện dưới ánh đèn an toàn mang lại trải nghiệm cảm xúc mà thế giới số không thể mang lại.
\end{itemize}

\subsection{Thực trạng phân mảnh của dịch vụ phòng tối và thiết bị chuyên dụng tại Việt Nam}
Mặc dù nhu cầu sáng tác ảnh phim tăng trưởng theo cấp số nhân trong 5 năm trở lại đây, hạ tầng kỹ thuật và dịch vụ hậu cần phục vụ cho cộng đồng này tại các đô thị lớn của Việt Nam (Hà Nội, TP. Hồ Chí Minh, Đà Nẵng, Cần Thơ) lại đang bộc lộ những hạn chế và bất cập nghiêm trọng:

\begin{enumerate}[leftmargin=2em]
    \item \textbf{Chi phí đầu tư thiết bị cá nhân quá đắt đỏ và khan hiếm:} 
    Một phòng tối chuẩn mực đòi hỏi hàng loạt thiết bị đặc thù đã ngừng sản xuất hàng loạt từ lâu: máy phóng ảnh (Enlarger như Durst, Beseler, Meopta), ống kính phóng ảnh chuyên dụng, máy sấy bản in ảnh, máy quét phim phân giải cao. Chi phí để một cá nhân tự xây dựng phòng tối hoàn chỉnh có thể lên tới hàng chục đến hàng trăm triệu đồng, chưa kể diện tích phòng kín, hệ thống cấp thoát nước và thông gió bắt buộc.
    
    \item \textbf{Quản lý và tiêu hủy hóa chất độc hại thiếu chuẩn mực:}
    Quá trình tráng và rọi ảnh sử dụng các hóa chất như thuốc hiện hình, thuốc hãm và thuốc định hình chứa ion bạc kim loại nặng. Tại các cơ sở nhỏ lẻ tự phát, việc pha chế sai tỷ lệ dẫn đến hỏng phim của khách, đồng thời việc xả thải nước rửa chứa ion bạc tự do trực tiếp ra hệ thống thoát nước sinh hoạt gây ô nhiễm môi trường nghiêm trọng nếu không có quy trình thu gom và trung hòa sinh thái đạt chuẩn~\cite{rogers2023sustainable}.
    
    \item \textbf{Phương thức vận hành thủ công, phân mảnh và thiếu chuyên nghiệp:}
    Hầu hết các phòng tối độc lập hiện nay vẫn tiếp nhận lịch đặt phòng của người dùng thông qua tin nhắn cá nhân (Facebook, Zalo) hoặc ghi chép vào sổ tay truyền thống. Phương thức này gây ra hàng loạt vấn đề:
    \begin{itemize}
        \item \textit{Trùng lặp lịch đặt phòng (Overbooking):} Khi hai hoặc nhiều khách hàng cùng liên hệ hoặc xác nhận qua nhiều kênh khác nhau, chủ cơ sở rất dễ nhận trùng một khung giờ cho nhiều người.
        \item \textit{Tranh chấp thiết bị và vật tư:} Khách hàng đến nơi mới phát hiện máy quét hoặc máy phóng ảnh chuyên dụng đang có người khác sử dụng, hoặc thuốc hóa chất đã hết hạn mà chưa kịp bổ sung.
        \item \textit{Thiếu minh bạch về tài chính:} Không có hóa đơn điện tử, phụ thu thời gian sử dụng vượt giờ hoặc bồi thường hư hỏng thiết bị không có cơ sở dữ liệu đối soát rõ ràng.
    \end{itemize}
\end{enumerate}

\subsection{Tính cấp thiết của đề tài nghiên cứu}
Từ những phân tích thực trạng nêu trên, việc ứng dụng các giải pháp công nghệ thông tin tiên tiến nhằm số hóa, chuẩn hóa và tự động hóa quy trình quản lý hệ sinh thái dịch vụ nhiếp ảnh phim là một đòi hỏi vô cùng cấp thiết. Cần phải có một nền tảng chuyển đổi số chuyên nghiệp với cơ sở dữ liệu tập trung, hỗ trợ quản lý định danh người dùng đa vai trò, kiểm soát tài nguyên thời gian thực và đặc biệt là giải quyết triệt để bài toán xử lý đồng thời (\textit{Concurrency}) khi đặt lịch phòng tối và thuê thiết bị đi kèm~\cite{habibi2024towards, zhang2023adaptive}.

Đề tài \textbf{"Nền tảng kết nối cộng đồng nhiếp ảnh phim với các dịch vụ phòng tối và phòng chụp"} không chỉ mang ý nghĩa lý thuyết trong việc vận dụng các kỹ thuật quản trị cơ sở dữ liệu hiện đại (T-SQL, Stored Procedure, các cấp độ cô lập giao dịch, chỉ dẫn khóa dòng)~\cite{silberschatz2020database, microsoft2024locking} mà còn mang giá trị thực tiễn to lớn trong việc bảo tồn và phát triển một nét văn hóa nghệ thuật thị giác độc đáo trong kỷ nguyên số.
